\chapter{RAG 基础：检索增强问答系统的最小闭环}

\begin{introduction}
\item {\heiti 本章核心问题}：当答案不能只依赖模型记忆时，怎样为企业知识助手搭建一个最小但真正可用的 RAG 闭环，并知道问题先出在哪一段？
\item {\heiti 主案例推进到哪一步}：让企业知识助手第一次真正“读企业文档”，从通用聊天工具迈向有证据链的问答系统。
\item {\heiti 读完后能做什么}：能把 RAG 拆成离线建库和在线问答两段，知道文档该怎么入库、日志该怎么记、答案为什么必须带引用，以及检索失败和生成失败该怎么区分。
\end{introduction}

\section{学习目标}
\begin{itemize}
  \item 理解为什么企业问答系统通常需要 RAG，而不是只靠模型参数记忆。
  \item 掌握最小 RAG 闭环中的关键环节：文档准入、切分、索引、检索、拼接、生成、引用与拒答。
  \item 能为企业知识助手设计一条可复盘的离线建库与在线问答链路，而不是只堆一个向量库接口。
  \item 能在回答错误时，先判断是检索层、拼接层还是生成层出了问题。
\end{itemize}

\section{问题背景}
到目前为止，我们已经知道如何调用模型、如何控制输出、也知道何时该做微调。但企业场景里真正难的，往往不是“模型会不会说”，而是“模型说的依据到底来自哪里”。制度今天更新，流程下周变更，FAQ 每个月都在补充，很多高价值知识并不适合被写死进参数里，而更适合以文档、页面、表格和知识条目的形式被随时接入。

企业知识助手还有一个更强的约束：它不能只给一个“像是对的回答”，而必须能说明依据、指出来源、允许复核、接受追责。只要答案会影响审批、报销、权限、运维和客户沟通，系统就必须在“回答能力”之外，再具备“证据组织能力”。这正是 RAG 的工程价值。

因此，本章不把 RAG 当作一个潮流名词，而把它当作一条最小可运营链路来看待：先把知识变成可检索对象，再把检索结果变成可审计答案。只要这条链路建立起来，后续的 PDF 解析、召回增强、重排序和 GraphRAG 才有意义；反过来，如果这条基础链路都不稳，所有高级技巧都只是在坏输入上继续加复杂度。

\section{核心概念}
\subsection{什么是 RAG}
RAG（Retrieval-Augmented Generation）可以直接理解为“先取证，再作答”。它不是单一算法，也不是给模型外挂一个检索插件，而是一条系统链路：
\begin{enumerate}
  \item 把原始知识整理成可被索引的片段；
  \item 按统一规则记录这些片段的来源、位置和版本；
  \item 当用户提问时，从知识库里召回最相关的若干片段；
  \item 把问题、片段和回答约束一起送入生成模型；
  \item 输出带依据、可回放、必要时可拒答的答案。
\end{enumerate}

这个定义里最重要的不是“检索”和“生成”这两个词本身，而是二者之间的职责分离。没有 RAG 时，我们把“记住事实”和“组织语言”都压给模型参数；有了 RAG 之后，动态知识交给知识库层治理，语言表达交给生成模型处理。只要职责清晰，系统才能真正复盘。

\subsection{RAG 解决的是知识接入问题，不是所有问题}
很多团队一看到回答不准，就想把所有问题都交给 RAG。这个判断并不严谨。一个实用的分界线是：如果错误主要来自知识过期、企业术语、制度例外、跨文档查证或必须提供来源，那么问题更偏向知识接入，RAG 通常是优先手段。如果错误主要来自语气、格式、输出结构或多轮对话风格，那更像提示工程或微调问题。

因此，RAG 的适用场景不是“所有问答系统”，而是“答案需要依赖外部、可变、可追溯知识”的问答系统。企业知识助手就是典型代表：员工问“差旅住宿超标怎么处理”“VPN 权限申请需要哪些审批”“某产品功能是否支持私有化部署”，这些问题的正确答案都依赖最新文档，而不是依赖模型曾经学过什么。

\subsection{最小闭环中的离线段与在线段}
把 RAG 真的做成系统时，最好先把它拆成两段：

\textbf{离线段}负责知识准备，包括文档采集、清洗、切分、向量化、索引构建和元数据登记。它要回答的是：“系统之后有没有机会拿到正确证据？”

\textbf{在线段}负责问题响应，包括问题预处理、召回、拼接、生成、引用输出和日志记录。它要回答的是：“系统拿到候选证据后，能不能忠实地组织成可用答案？”

这个切分非常重要，因为很多团队会把所有问题都归结为“模型不够强”。实际上，离线段一旦切错，在线段不管再怎么调参，也只能在坏证据上作答；反过来，离线段做得不错但在线段拼接混乱，系统也会表现得像检索失败一样。工程上越早把两段职责拆清楚，后面的优化路径越不会走偏。

\begin{figure}[H]
\centering
\begin{tikzpicture}[
  font=\small,
  box/.style={llm box, minimum width=2.5cm, minimum height=1cm},
  arr/.style={llm arr}
]
\node[box] (docs) {文档准入};
\node[box, right=0.7cm of docs] (offline) {切分与建库};
\node[box, right=0.7cm of offline] (retrieve) {在线检索};
\node[box, right=0.7cm of retrieve] (prompt) {上下文拼接};
\node[box, right=0.7cm of prompt] (answer) {带引用作答};
\draw[arr] (docs) -- (offline);
\draw[arr] (offline) -- (retrieve);
\draw[arr] (retrieve) -- (prompt);
\draw[arr] (prompt) -- (answer);
\end{tikzpicture}
\caption{企业知识助手的最小 RAG 闭环}
\label{fig:minimal-rag-loop}
\end{figure}

\begin{table}[H]
\centering
\small
\begin{tabularx}{\textwidth}{>{\raggedright\arraybackslash}p{2.8cm}>{\raggedright\arraybackslash}p{3.2cm}>{\raggedright\arraybackslash}X}
\toprule
\textbf{环节} & \textbf{最小设计重点} & \textbf{最容易暴露的问题} \\
\midrule
文档准入 & 只纳入有效、授权、可追踪的知识源 & 过期制度混入、重复文档入库、权限边界不清 \\
切分与建库 & 保住语义边界、来源位置和版本信息 & 标题丢失、表格断裂、片段回溯不到原文 \\
在线检索 & 让相关证据稳定进入候选集合 & 召回不到关键条款、相似但无用的噪声过多 \\
上下文拼接 & 明确材料、问题、引用规则与拒答规则 & 片段顺序混乱、材料没真正进入模型上下文 \\
答案生成 & 忠于证据、证据不足时保守回答 & 无依据扩写、忽略例外条款、拒答不稳定 \\
\bottomrule
\end{tabularx}
\caption{最小 RAG 闭环中每一段真正要守住的东西}
\label{tab:rag-min-loop}
\end{table}

\subsection{文档准入先决定上限}
刚开始做 RAG 时，大家最容易把注意力放在向量库和模型上，却忽略了更上游的一个问题：哪些文档应该进库。企业知识助手不是“把所有文件夹全部导入”就算完成了，而是要先做一轮知识准入。

这一步至少要回答四个问题。第一，文档是否仍然有效，还是已经被新版本替代。第二，文档是否属于当前业务域，还是只是历史资料和低价值附件。第三，文档是否允许当前助手访问，尤其是包含审批、权限、人事或报价信息时。第四，文档能不能稳定定位来源，例如是否有文档编号、页面、章节或更新时间。

也就是说，最小 RAG 的起点不是索引，而是\textbf{知识治理}。如果准入阶段把无效材料和高噪声材料都放进来，后面的检索器只会更努力地把错误知识送给模型。

\subsection{文档切分不是技术细节，而是语义边界设计}
最小闭环里最容易被低估的一步，是切分。一个片段切得过长，会把无关信息一起带进检索；切得过短，又会把表头、定义、适用范围和例外条件拆散。尤其在制度型文档中，真正重要的往往不是一句结论，而是“适用对象 + 限制条件 + 例外条款”这组三件套。

因此，切分不该只盯着字符长度，而要先看语义边界。工程上通常有三种常见起点：
\begin{itemize}
  \item \textbf{固定长度切分}：实现简单，适合快速验证闭环，但容易切坏结构。
  \item \textbf{按标题和段落切分}：更符合文档阅读逻辑，适合制度、手册和 FAQ。
  \item \textbf{按结构块切分}：把表格、流程、代码块、告警示例等单独处理，适合复杂文档。
\end{itemize}

真正稳定的企业知识助手，往往不是三选一，而是“先识别结构，再对不同结构采用不同切分规则”。不过在最小闭环阶段，不需要一开始就追求最复杂方案，只要做到“片段语义基本完整，且能回到原文位置”，就已经比盲目按字符截断强很多。

\begin{table}[H]
\centering
\small
\begin{tabularx}{\textwidth}{>{\raggedright\arraybackslash}p{2.8cm}>{\raggedright\arraybackslash}p{2.8cm}>{\raggedright\arraybackslash}X}
\toprule
\textbf{文档形态} & \textbf{初始切分策略} & \textbf{原因} \\
\midrule
制度、手册、规范 & 按标题与段落切分 & 保留章节层级和适用范围，便于引用 \\
FAQ、工单总结 & 一问一答成块 & 问题与答案必须绑定，不能拆散 \\
表格、报价单、配置表 & 结构块切分 & 表头、列名和单元格关系比长度均匀更重要 \\
流程文档、操作步骤 & 按步骤块切分 & 顺序信息属于语义的一部分，不能打断 \\
\bottomrule
\end{tabularx}
\caption{企业知识助手中常见文档的初始切分策略}
\label{tab:rag-chunking-start}
\end{table}

\subsection{元数据不是附属品，而是可追溯性的骨架}
一个能上线的 RAG 片段，不应该只是一段裸文本。至少还应带上文档编号、标题路径、页码或段落编号、更新时间、权限域、来源链接和版本号中的若干项。没有这些元数据，系统即使偶尔答对了，也很难真正复盘。

元数据的价值主要体现在三个方面。第一，它让答案可以输出引用，而不是只给一个抽象结论。第二，它让团队能在出错后快速定位“模型到底看了什么”。第三，它让知识库更新具备可治理性，例如制度更新后可以按版本回滚、重建或差分替换。

对企业知识助手来说，最有用的一组最小元数据通常包括：\texttt{document\_id}、\texttt{title\_path}、\texttt{page}、\texttt{updated\_at}、\texttt{owner}、\texttt{permission\_scope}。这组字段不一定一次配齐，但至少要知道自己将来要往这个方向补。

\subsection{权限过滤要发生在召回前后两次，而不是等答案写完再补救}
很多团队第一次做企业 RAG 时，会把权限问题理解成“最后把答案页面藏住就行”。这远远不够。因为一旦越权片段先被召回并进入上下文，它就已经污染了答案生成过程；即使最后不展示原文，模型也可能已经把其中的结论写进回答。

因此，更稳的做法通常是做两次权限控制。第一次发生在\textbf{召回前或召回时}，利用权限域字段、部门域、文档级 ACL 或知识库分桶，把当前用户本就不该看到的材料排除在候选集合之外。第二次发生在\textbf{生成前或生成后}，再检查最终入模片段和引用结果，确认没有把跨域材料混进来。

\begin{table}[H]
\centering
\small
\begin{tabularx}{\textwidth}{>{\raggedright\arraybackslash}p{2.8cm}>{\raggedright\arraybackslash}p{3.2cm}>{\raggedright\arraybackslash}X}
\toprule
\textbf{权限控制位置} & \textbf{主要作用} & \textbf{为什么不能只做这一层} \\
\midrule
召回前 / 召回时过滤 & 从候选集合里直接排除越权材料 & 这是最关键的一层，但仍可能因为配置错误或索引污染漏进少量材料 \\
入模前二次校验 & 防止拼接阶段把跨域片段带进 prompt & 召回正确不代表拼接一定安全，尤其在多源混合检索时 \\
答案输出后展示控制 & 约束最终可见引用与原文跳转 & 只能兜底展示，无法逆转模型已经读过越权材料这件事 \\
\bottomrule
\end{tabularx}
\caption{企业知识助手里的权限控制最好分布在召回、入模和输出三个位置}
\label{tab:rag-permission-guard}
\end{table}

对主案例来说，这个设计尤其关键。比如“权限申请说明”与“高权限审批例外名单”可能都与用户问题高度相关，但它们的可见范围并不相同。系统如果没有把权限过滤做进召回层，就很容易出现“答案似乎合理，却引用了不该给员工看的内部规则”这种高风险错误。

\subsection{Embedding 在 RAG 里到底做什么}
上一章已经讲过 embedding 的基本直觉，这里把它落回 RAG。对检索器来说，embedding 的工作不是“把文本变高级”，而是把问题和文档片段投到同一个可比较的表示空间里。之后系统再通过点积、余弦相似度或近似最近邻搜索，判断哪些片段与当前问题更接近。

这背后至少有三层工程判断。第一，\textbf{什么对象被编码}。是整段正文、标题加正文、FAQ 问句、表格行，还是摘要侧车索引，它们会直接影响检索行为。第二，\textbf{问题和文档是否真的在同一种表达空间里}。若用户用口语提问，文档却全是制度标题和编号，embedding 再强也可能先天错位。第三，\textbf{领域语言是不是被覆盖}。产品名、内部系统名、报错码和制度编号越多，通用 embedding 的“常识语义”就越不够。

因此，最小闭环阶段真正需要守住的，不是“先选一个最热门 embedding 模型”，而是确认它至少满足三件事：中文或多语种支持可靠、短字段与术语检索不至于完全失真、问题和文档能在同一空间里形成稳定邻近关系。只要这三件事没守住，后面的混合检索、重排序甚至 SFT 都很难真正救回来。

\subsection{近似最近邻与元数据过滤：最小系统也要懂的检索底盘}
很多书会把向量库讲成一个黑箱：把向量扔进去，再把近邻取出来。工程上更实用的理解是，向量检索通常分成两层。第一层是\textbf{近似最近邻搜索}，负责在大规模片段集合里快速找到一批大致相关的候选；第二层是\textbf{元数据过滤}，负责把这些候选限制在当前业务域、时间范围、权限范围或文档类型里。

这两层缺一不可。只有近似最近邻，没有元数据过滤时，系统很容易把“语义像但业务域不对”的材料也召回来。只有过滤，没有语义检索时，系统又可能在正确范围里找不到真正相关的条款。最小闭环不要求一开始就把索引调到极致，但至少要知道：检索并不只是“相似度排序”，它同时也是“候选集合裁剪”。

\begin{note}
很多企业问答的第一性问题不是“向量库够不够快”，而是“候选集合一开始有没有被裁对”。只要部门域、产品线、版本范围或文档类型没有先限住，后面再精细的排序都会在大噪声池里挣扎。
\end{note}

\subsection{最小检索策略要先稳，再谈花样}
一提到 RAG，很多人马上就会想到“向量检索”。实际上，最小闭环阶段的目标不是用上最多检索技术，而是让相关证据稳定出现在候选集合里。对不少企业知识场景来说，如果文档里有大量术语、编号、命令名和产品名，只做语义检索未必稳，关键词检索反而更重要。

因此，最小检索策略可以这样理解：
\begin{itemize}
  \item 如果问题与文档表述高度接近，单一路径的向量检索或关键词检索都可能足够。
  \item 如果问题中包含制度编号、配置项、命令名、接口名，关键词能力不能缺席。
  \item 如果用户表达口语化，而文档语言较正式，后续才考虑查询改写或混合检索。
\end{itemize}

也就是说，最小闭环不要求一开始就把混合检索、重排序、多查询全部上齐；它要求的是你知道自己当前的检索器为什么选它，以及它最可能漏掉什么。

\subsection{关键词、向量与混合检索：最小系统也要知道各自短板}
旧稿里反复强调过一个很实用的判断：{\heiti 企业知识场景里的检索，从来不是“向量检索一统天下”。} 因为很多高价值线索并不是抽象语义，而是具体字段、命令名、制度编号、接口名、报错码和版本号。

\begin{table}[H]
\centering
\small
\begin{tabularx}{\textwidth}{>{\raggedright\arraybackslash}p{2.6cm}>{\raggedright\arraybackslash}p{3cm}>{\raggedright\arraybackslash}p{2.8cm}>{\raggedright\arraybackslash}X}
\toprule
\textbf{检索路线} & \textbf{更擅长什么} & \textbf{最容易漏什么} & \textbf{最小闭环中的建议} \\
\midrule
关键词 / BM25 & 术语、编号、命令、报错码、精确字段 & 口语问法、同义改写、隐式语义 & 术语密集场景优先保留，不能因为“老”就删掉 \\
向量检索 & 口语问法、语义近似、表达改写 & 编号、短字段、精确字面线索 & 适合作为通用召回主路，但要知道它在哪些场景发虚 \\
混合检索 & 同时兼顾语义近似和精确关键词 & 配置复杂度更高，需要做融合规则 & 当文档既有术语又有自然语言说明时，往往是最稳的默认解 \\
\bottomrule
\end{tabularx}
\caption{最小 RAG 也应知道关键词检索、向量检索和混合检索的职责边界}
\label{tab:retrieval-route-choice}
\end{table}

对企业知识助手来说，一个很典型的例子是“VPN-4037 报错怎么处理”和“差旅制度 2025-04 版是否允许高铁改签”。这类问题里，编号本身就是关键信号。若检索器只擅长“语义相似”，却抓不住这些精确线索，系统就会在看起来最简单的问题上翻车。

\subsection{查询改写不是默认开关，而是症状驱动修复}
旧稿里讲过查询改写、多查询检索和伪文档生成。这些方法在 RAG 里确实有价值，但放到最小闭环阶段，最重要的不是“先上最全的一套”，而是先搞清楚\textbf{什么时候它们真的能修问题，什么时候只是在增加链路噪声。}

查询改写最适合处理三类症状。第一，用户问题过于口语、省略或上下文依赖强，例如“那高铁改签呢”“这个例外怎么报”。第二，文档语言很正式，而用户问法很口语化，二者在字面上几乎碰不到。第三，用户问题其实同时包含多个子问题，需要先拆开才能更稳定地命中文档。除此之外，若问题本身已经很清楚，检索仍然失败，往往更该先回头检查知识覆盖、切分和索引，而不是默认多加一层 LLM 改写。

\begin{table}[H]
\centering
\small
\begin{tabularx}{\textwidth}{>{\raggedright\arraybackslash}p{3cm}>{\raggedright\arraybackslash}p{3.2cm}>{\raggedright\arraybackslash}X}
\toprule
\textbf{观察到的症状} & \textbf{更值得尝试的改写动作} & \textbf{为什么先做这个} \\
\midrule
问题过短、强依赖上一轮历史 & 先做会话感知改写，把省略问法补成完整检索问句 & 让检索器看到明确对象，而不是一串上下文省略词 \\
用户口语与文档书面语差异很大 & 做术语补齐或规范化改写 & 提高问题表达与文档表达的重合度 \\
一个问题里同时混着多个子任务 & 先拆子查询，再合并证据 & 避免单条查询在多个意图之间平均失焦 \\
问题已经清楚但仍检索不到 & 暂不优先改写，先查知识覆盖与切分 & 避免用额外 LLM 调用掩盖底层知识问题 \\
\bottomrule
\end{tabularx}
\caption{查询改写应由具体症状触发，而不是默认插在所有 RAG 链路中}
\label{tab:rag-query-rewrite-when}
\end{table}

这也是为什么第 8 章才会继续展开 HyDE、RAG-Fusion 等增强策略。最小闭环阶段先要建立的是判断能力：\textbf{当前失败是因为“问法难检索”，还是因为“知识本身没进来”。} 只有把这件事分清，多查询和改写才会成为增益，而不是额外噪声源。

\subsection{二级索引不是花样，而是把“好检索”和“好生成”拆开}
旧稿里还讲过一种很有价值的思路：不要让“检索用的表示”和“生成用的原文”永远绑成同一块。很多时候，更适合检索的是标题、关键句、摘要或结构化提要；而真正适合交给大模型生成答案的，则是对应的原始段落、表格或完整条款。

这就形成了一个很实用的二级索引思路：
\begin{itemize}
  \item \textbf{第一级索引}：用更利于召回的短表示参与检索，比如标题、摘要、FAQ 问句、关键字段。
  \item \textbf{第二级映射}：把召回结果映射回原始文本块、表格块或完整条文，再交给模型阅读。
\end{itemize}

它的价值不在“更高级”，而在于解决一个常见矛盾：\textbf{适合被 embedding 的，不一定是最适合直接进 prompt 的；适合直接进 prompt 的，又不一定最容易被召回。} 对最小闭环来说，你不一定一开始就上完整二级索引系统，但至少要有这个意识。否则团队会不断在“把块切短一点方便检索”和“把块保完整一点方便生成”之间来回摇摆。

\subsection{\texorpdfstring{$top$-$k$}{top-k} 不是越大越好}
很多团队在召回不准时，会本能地把 \texttt{top-k} 越调越大，觉得“给模型更多材料，总不会更差”。这往往会带来相反效果。召回越多，无关材料越多，提示越长，模型越难抓住真正关键的证据。对于需要精确查条款的企业问答，过长上下文反而会让答案变得模糊。

因此，\texttt{top-k} 不该被当作越大越保险的旋钮，而应该被当作“证据密度”和“上下文预算”之间的平衡器。一个可操作的做法，是根据问题类型先给出初始范围，而不是全场景统一一个数值。

\begin{table}[H]
\centering
\small
\begin{tabularx}{\textwidth}{>{\raggedright\arraybackslash}p{3cm}>{\centering\arraybackslash}p{1.8cm}>{\raggedright\arraybackslash}X}
\toprule
\textbf{问题类型} & \textbf{初始 \texttt{top-k}} & \textbf{说明} \\
\midrule
单条制度查证 & 3--5 & 重点是把最直接的条款拿准，噪声越少越好 \\
FAQ 类问答 & 4--6 & 往往需要一个主答案和一两个补充说明 \\
跨文档综合问题 & 6--8 & 需要多个来源，但仍要控制上下文拥挤 \\
需要表格与正文联合解释的问题 & 5--8 & 既要召回条文，也要召回表格或附件 \\
\bottomrule
\end{tabularx}
\caption{企业知识助手中常见问题的初始 \texttt{top-k} 建议}
\label{tab:rag-topk-start}
\end{table}

\subsection{候选去重与多样性控制，往往比继续增大 \texttt{top-k} 更值钱}
召回阶段还有一个非常常见却常被忽略的问题：\textbf{系统看上去召回了很多片段，实际上大半是同一页、同一条款或同一段话的不同切片。} 这种“表面上很多、实质上很重复”的候选集合，会同时伤害两个目标。一方面，它挤占上下文预算，让真正互补的证据进不来；另一方面，它会给模型一种假象，仿佛某一条局部措辞被多份材料重复确认了。

因此，在最小 RAG 闭环里，候选去重和多样性控制通常比继续增大 \texttt{top-k} 更先值得做。所谓多样性，不是故意追求“各种不同风格的段落”，而是确保候选里能覆盖\textbf{主条款、例外条件、附件说明、表格字段}这些互补信息，而不是让同一条主条款占满前几位。

\begin{table}[H]
\centering
\small
\begin{tabularx}{\textwidth}{>{\raggedright\arraybackslash}p{2.8cm}>{\raggedright\arraybackslash}p{3cm}>{\raggedright\arraybackslash}X}
\toprule
\textbf{候选集合问题} & \textbf{最直接后果} & \textbf{更稳的第一步修法} \\
\midrule
同一文档相邻切片过多 & 互补证据进不来，上下文变窄 & 先按文档或父块分桶，再限制每桶入选数 \\
同义重复片段过多 & 模型误以为单一结论被多次独立支持 & 先做近重复去重，再看是否还需要增大 \(k\) \\
只有主条款，没有例外条件 & 回答容易过度肯定 & 对高风险问题显式保留“例外 / 附件 / 备注”类候选 \\
只有正文，没有表格或附件 & 答案缺字段、缺阈值、缺条件 & 为不同文档形态保留最低覆盖比例 \\
\bottomrule
\end{tabularx}
\caption{RAG 候选集合的质量，不只取决于数量，也取决于互补性}
\label{tab:rag-candidate-diversity}
\end{table}

\subsection{提示拼接要把规则写清楚}
检索结果并不会自动变成好答案。系统还必须清楚地告诉模型：哪些是材料，哪些是用户问题，回答是否必须引用，材料不足时该怎么做，材料互相矛盾时是否要提示人工确认。这些规则如果不显式写进提示词，模型很容易把检索到的材料当作“可参考但不必严格遵守”的背景信息。

一个可上线的企业知识助手，至少要有三条回答约束：
\begin{itemize}
  \item 仅基于提供的材料作答，不从常识自行扩写结论；
  \item 给出答案时标注来源片段或文档位置；
  \item 当证据不足、证据冲突或权限不足时，允许拒答或提示人工确认。
\end{itemize}

这并不是在“限制模型能力”，而是在明确系统责任。对企业问答来说，保守比自信更重要，可回溯比华丽措辞更重要。

\begin{lstlisting}[language=Python]
retrieved_docs = retriever.invoke(question)
selected_docs = retrieved_docs[:4]

context = "\n\n".join(
    f"[{doc.metadata['document_id']}|p.{doc.metadata['page']}] {doc.page_content}"
    for doc in selected_docs
)

prompt = f"""
你是企业知识助手。请仅依据给定材料回答问题。
如果材料不足，请直接回答“依据不足，建议人工确认”。
回答时请引用材料编号。

材料：
{context}

问题：
{question}
"""
\end{lstlisting}

上面的代码并不复杂，但它把三个关键接口都暴露出来了：检索到了什么、最终送进模型的是什么、答案是否被约束为带引用且可拒答。只要这三个接口可见，团队就具备继续优化的基础。

\subsection{回答契约要固定：结论、依据、风险提示按同一顺序出现}
企业知识助手并不只是“把材料交给模型”就结束了。对很多高价值问题，更关键的是回答的组织顺序要稳定。一个最实用的做法，是把回答契约固定成三段：先给结论，再给依据，最后给风险提示或例外条件。这样做的好处是，用户和审核者都会知道该从哪里看“最终判断”，从哪里看“证据支持”，又从哪里看“还有哪些边界条件”。

如果不固定这个顺序，系统很容易出现三种不稳定行为。第一，先铺材料后给结论，导致用户误读。第二，只给结论不解释依据，失去审计性。第三，结论看似明确，却把“需要额外审批”“仅适用于某岗位”“材料不足时应升级人工”这些风险提示藏在段落末尾，等于把最关键的约束弱化了。

\begin{table}[H]
\centering
\small
\begin{tabularx}{\textwidth}{>{\raggedright\arraybackslash}p{2.6cm}>{\raggedright\arraybackslash}p{3cm}>{\raggedright\arraybackslash}X}
\toprule
\textbf{回答段落} & \textbf{最小要求} & \textbf{企业知识助手中的意义} \\
\midrule
结论 & 直接回答当前问题，不绕弯 & 方便员工先获得可执行判断 \\
依据 & 列出条款、页码、表格行或文档编号 & 方便复核，也方便后续追责与纠错 \\
风险提示 / 例外条件 & 明确限制范围、补充材料和升级人工条件 & 防止系统把“通常情况”误包装成“绝对规则” \\
\bottomrule
\end{tabularx}
\caption{企业知识助手中的最小回答契约应把结论、依据和风险提示拆开}
\label{tab:rag-answer-contract}
\end{table}

这类契约看上去像输出格式问题，实际上会反过来影响检索与日志设计。因为一旦系统承诺“必须给依据、必须说明风险提示”，团队就会更自然地去保留元数据、记录引用链、补齐例外条款片段，而不是只盯着生成流畅度。

\subsection{证据冲突时先显式标注冲突，不要替制度拍板}
企业知识场景里还有一类最容易被低估的问题：\textbf{召回到的证据并不是全都一致。} 新旧制度并存、附件表格和正文口径不一、不同部门文档对同一流程有不同表述，这些都很常见。若系统在这种情况下仍然被要求“给一个确定答案”，模型往往会自己替组织做拍板，而这恰恰是最危险的行为。

因此，最小闭环里也应该把“冲突证据如何处理”写成协议。更稳妥的顺序通常是：先指出冲突来自哪两份材料，再说明当前无法仅凭现有证据做最终判断，最后给出补查方向或升级人工动作。这样做的目标不是让系统显得保守，而是让它保持\textbf{证据透明和责任边界透明。}

\begin{table}[H]
\centering
\small
\begin{tabularx}{\textwidth}{>{\raggedright\arraybackslash}p{2.8cm}>{\raggedright\arraybackslash}p{3cm}>{\raggedright\arraybackslash}X}
\toprule
\textbf{冲突类型} & \textbf{更稳的回答策略} & \textbf{为什么不能直接求平均} \\
\midrule
新旧版本条款不一致 & 优先报告版本冲突，并提示确认现行版本 & “折中”答案通常在制度上没有合法性 \\
正文与附件表格口径不同 & 同时引用两处，并要求补查说明或人工确认 & 模型无法替代制度 owner 裁决哪处生效 \\
不同部门文档边界不同 & 先说明适用域差异，再提示按部门或权限确认 & 不同组织域本来就可能没有统一答案 \\
\bottomrule
\end{tabularx}
\caption{证据冲突时，系统应先暴露冲突，而不是替组织完成裁决}
\label{tab:rag-conflict-handling}
\end{table}

\subsection{RAG 不是 SFT 的替身：知识接入和行为塑形要分开看}
很多团队在做到这里时会追问：既然目标是让企业知识助手答好制度问题，为什么不直接做 SFT？更稳妥的回答是：\textbf{RAG 和 SFT 注入的是两种不同东西。} RAG 注入的是外部、可变、可追溯的知识；SFT 或 PEFT 塑造的是回答风格、输出格式、拒答规则和任务习惯。

如果问题的核心矛盾是“答案依赖最新制度、附件表格、权限边界和可引用证据”，那么优先级通常应放在 RAG，因为这些东西本来就不该只靠参数记忆保存。相反，如果问题的核心矛盾是“材料明明够，但回答总是不按模板写、不稳定给引用、不按风险要求拒答”，那就更像是任务行为问题，适合用提示工程、SFT 或 PEFT 去塑形。

\begin{table}[H]
\centering
\small
\begin{tabularx}{\textwidth}{>{\raggedright\arraybackslash}p{2.8cm}>{\raggedright\arraybackslash}p{3.2cm}>{\raggedright\arraybackslash}X}
\toprule
\textbf{手段} & \textbf{更适合解决什么} & \textbf{企业知识助手中的典型场景} \\
\midrule
RAG & 最新知识接入、证据引用、权限边界、可回溯事实 & 制度问答、产品 FAQ、表格条款查证、需要给出处的回答 \\
SFT / PEFT & 输出格式、拒答风格、术语口径、任务流程稳定性 & 固定模板答复、结构化归档、统一引用格式、保守拒答规范 \\
RAG + SFT & 同时需要最新知识和稳定行为 & 有依据地回答制度问题，并按统一格式给结论、引用和风险提示 \\
\bottomrule
\end{tabularx}
\caption{RAG 与 SFT 不是二选一，而是分别解决知识接入和行为塑形}
\label{tab:rag-vs-sft}
\end{table}

对这本书的主案例来说，更稳的顺序通常是：先让最小 RAG 闭环跑通，确保系统能稳定拿到证据；再用提示工程或小规模 SFT 把“如何基于证据说话”这件事做稳。反过来，如果一开始就试图把整套制度知识全塞进 SFT，一方面更新代价高，另一方面一旦制度换版，模型参数里的旧知识也很难治理。

\subsection{日志和回放能力决定系统能否复盘}
很多团队上线一个 RAG 原型后，最先遇到的问题不是效果差，而是效果差的时候看不到证据链。用户只说“答错了”，团队却不知道是检索错、拼接错还是生成错。要避免这种状态，系统必须记录最小回放日志。

一个最有用的最小日志集通常包括：原始问题、归一化问题、召回的候选片段、最终入模的片段、提示模板版本、模型版本、最终答案和引用来源。没有这些字段，就无法把一次错误回答重新跑回来。

\begin{table}[H]
\centering
\small
\begin{tabularx}{\textwidth}{>{\raggedright\arraybackslash}p{3cm}>{\raggedright\arraybackslash}p{3cm}>{\raggedright\arraybackslash}X}
\toprule
\textbf{日志字段} & \textbf{用途} & \textbf{为什么不能省} \\
\midrule
原始问题与归一化问题 & 复盘用户到底问了什么 & 口语化问题经过改写后，检索路径可能完全不同 \\
召回候选列表 & 判断检索器是否找到证据 & 看不到候选，就无法区分检索失败与生成失败 \\
最终入模片段 & 判断拼接阶段是否丢材料 & 候选正确不代表真正进了模型上下文 \\
提示模板版本 & 判断规则是否改变了行为 & 很多回归来自提示改版而非模型变差 \\
模型版本与答案 & 识别生成层差异 & 同样的证据在不同模型上可能会被不同解读 \\
\bottomrule
\end{tabularx}
\caption{RAG 系统最值得保留的最小回放日志}
\label{tab:rag-logs}
\end{table}

\subsection{知识库不是一次建完：最小闭环也要有更新纪律}
很多原型系统的隐性问题，不是第一次建库建得差，而是建完以后再也不更新。企业知识助手一旦进入真实环境，就会面对制度换版、流程更新、FAQ 变化、权限边界调整。此时如果索引仍然停在旧版本，RAG 就会稳定地产出“带引用的旧答案”。

因此，最小闭环阶段最好就建立三条更新纪律：
\begin{itemize}
  \item \textbf{文档入库要带版本和时间戳}：否则新旧制度混在一起时很难裁决。
  \item \textbf{更新策略要明确是增量替换还是全量重建}：不同知识类型的治理方式不一样。
  \item \textbf{高风险知识要有 owner}：谁负责确认制度是否已过期，不能让向量库自己承担治理职责。
\end{itemize}

这也是为什么前面一直强调元数据。没有版本、更新时间和责任人字段，知识库看起来是“可检索”的，但并不是“可运营”的。

\subsection{增量替换还是全量重建：更新策略要按知识类型分开}
旧稿里关于建库的一个经验很值得保留：不是所有知识更新都适合同一套刷新方式。FAQ、新增知识条目和操作手册的小修订，往往适合做增量替换；制度全集、标题结构整体变化或表格模板改版，则常常更适合全量重建。原因很简单，前者的语义边界和引用关系变化较局部，后者则可能把切分、标题路径和表格绑定关系整体改写。

\begin{table}[H]
\centering
\small
\begin{tabularx}{\textwidth}{>{\raggedright\arraybackslash}p{3cm}>{\raggedright\arraybackslash}p{2.8cm}>{\raggedright\arraybackslash}X}
\toprule
\textbf{知识变化类型} & \textbf{更常用的刷新方式} & \textbf{为什么} \\
\midrule
新增 FAQ 或知识条目 & 增量写入 & 新片段边界清楚，对旧索引破坏小 \\
小范围条款修订 & 定位后替换相关片段 & 能保留大部分稳定索引，只改动受影响区域 \\
制度大版本换版 & 全量重建更稳 & 标题路径、页码、条款顺序和引用关系可能整体变化 \\
表格模板改版 & 常需全量重建或整表替换 & 行列关系一旦变化，旧切分和旧引用往往一起失效 \\
\bottomrule
\end{tabularx}
\caption{RAG 知识刷新策略应跟随知识变化类型，而不是一刀切}
\label{tab:rag-refresh-strategy}
\end{table}

这条纪律对主案例很关键。像“员工差旅制度”这类高风险知识，如果只是把新 PDF 直接追加进库而不淘汰旧版本，系统就会在同一问题上同时召回新旧条款，看起来像模型“前后矛盾”，其实是知识层已经失序。

\subsection{先判断是哪一层失败}
RAG 最实用的工程直觉之一，就是在回答错误时先分层归因。一个简单判断方法如下：
\begin{itemize}
  \item 如果召回片段本身就与问题无关，优先判为\textbf{检索失败}。
  \item 如果召回片段相关，但没有真正进入最终提示，优先判为\textbf{拼接失败}。
  \item 如果召回片段相关且已经入模，但答案仍不忠于材料，优先判为\textbf{生成失败}。
  \item 如果答案引用的是旧版或越权材料，优先判为\textbf{知识治理失败}。
\end{itemize}

这个分层判断看似朴素，却决定了整个优化方向。只要团队习惯先问“哪一层错了”，而不是先说“模型又不行”，后续的调优成本就会大幅下降。

\begin{table}[H]
\centering
\small
\begin{tabularx}{\textwidth}{>{\raggedright\arraybackslash}p{3cm}>{\raggedright\arraybackslash}p{3cm}>{\raggedright\arraybackslash}X}
\toprule
\textbf{观察到的现象} & \textbf{优先怀疑层} & \textbf{第一动作} \\
\midrule
召回片段和问题明显无关 & 检索层 & 先查切分、索引方式、查询表达和 \texttt{top-k} \\
候选片段相关，但最终提示里没有 & 拼接层 & 先查上下文预算、去重逻辑和模板拼接顺序 \\
提示里已有正确证据，但答案仍扩写 & 生成层 & 先收紧回答约束，检查引用和拒答规则 \\
答案引用了旧制度或越权信息 & 知识治理层 & 先查文档准入、版本淘汰和权限过滤 \\
\bottomrule
\end{tabularx}
\caption{从现象反推最小 RAG 闭环中的故障位置}
\label{tab:rag-routing}
\end{table}

\section{主案例推进}
在企业知识助手的主案例里，我们把第一阶段目标收得很明确：不追求“回答一切”，而是先把报销制度、VPN 权限申请说明、产品 FAQ 和运维操作手册四类高价值文档做成一个可复盘的最小闭环。

\subsection{先从高价值文档开始，而不是全量入库}
项目起步时，我们不把共享盘所有文件都接进系统，而是先选一批高频、高争议、更新可控的文档。这样做有三个好处。第一，问题空间收窄，便于快速验证。第二，噪声更少，更容易看出检索器是否真的有效。第三，后续复盘时，知识治理成本可控。

在这个案例里，第一批入库材料是：
\begin{enumerate}
  \item 最新版员工差旅与报销制度；
  \item VPN、代码仓库、工单系统等权限申请说明；
  \item 产品功能 FAQ 与私有化部署说明；
  \item 运维操作手册中的高频排障流程。
\end{enumerate}

这些材料有一个共同特征：它们都直接影响员工日常工作，而且答案通常需要给出明确依据。如果把这批材料做稳，系统就已经从“会聊天”迈向“能支撑真实业务”。

\subsection{离线建库要先产出一套固定清单}
为了让团队对离线段有共同语言，我们把建库流程固定成下面这张清单：
\begin{enumerate}
  \item 记录文档来源、负责人、更新时间和版本号；
  \item 去除失效文档与重复附件；
  \item 按标题、段落、表格等结构做切分；
  \item 为每个片段补齐文档编号、章节路径、页码和权限域；
  \item 建立初始索引，并抽样检查若干问题的召回结果；
  \item 保存一份可回放的建库配置，包括切分参数与索引版本。
\end{enumerate}

这套清单的意义不只是“把文档导进去”，而是让知识库层第一次具备可维护性。以后制度更新时，团队知道要替换哪一批片段；权限边界变化时，也知道该过滤哪一类材料。

\subsection{在线问答要让证据链完整可见}
假设员工提问：“出差住宿超标后还能报销吗？”系统的标准链路应该长这样：
\begin{enumerate}
  \item 接收原始问题，并保留原句进入日志；
  \item 检索报销制度中与“住宿标准”“超标”“审批例外”相关的若干片段；
  \item 选出真正相关的 3--4 个片段，连同文档编号和页码送入提示模板；
  \item 要求模型先给结论，再给依据，再说明是否存在例外条款；
  \item 若证据仅说明“标准金额”，却没有说明“超标后是否允许补差或审批”，则输出“依据不足，建议人工确认或补充例外审批制度”。
\end{enumerate}

这条链路的重点不在于模型回答得多流畅，而在于每一步都能看到。只要“问题-候选-入模-答案”这四段是透明的，团队就知道哪里还有提升空间。

\begin{lstlisting}[language=Python]
trace = {
    "question": "出差住宿超标后还能报销吗？",
    "retrieved": [
        "TRAVEL-2026 p.12 住宿标准",
        "TRAVEL-2026 p.14 超标审批例外",
        "FIN-FAQ p.3 报销附件要求"
    ],
    "selected": [
        "TRAVEL-2026 p.12 住宿标准",
        "TRAVEL-2026 p.14 超标审批例外"
    ],
    "prompt_version": "rag_answer_v3",
    "answer": "可以申请，但需满足制度中的超标审批条件，并附审批记录。"
}
\end{lstlisting}

这个 \texttt{trace} 看起来像一段普通日志，但它其实对应了整个最小闭环的核心资产。团队以后要做评测、回归测试、错误归因和知识更新，都会反复依赖这种最小回放能力。

\subsection{最小闭环真正要验证的不是“能答”，而是“能稳定答”}
如果只是演示给同事看，一个只在少数样例上表现不错的原型就够了。但要把企业知识助手继续往后推进，我们验证的重点应该变成三件事：
\begin{itemize}
  \item 同一类问题多次提问时，是否能稳定召回到同类证据；
  \item 当文档没有答案时，是否会稳定拒答而不是自信编造；
  \item 当制度更新后，答案是否会切换到新版本，而不是继续引用旧材料。
\end{itemize}

也就是说，RAG 的最小闭环阶段，真正重要的是\textbf{稳定性}而不是\textbf{偶然的高光答案}。只有这一步稳住，下一章的文档解析、召回增强和 GraphRAG 才值得引入。

\section{常见坑与工程取舍}
\begin{itemize}
  \item \textbf{一开始就追求复杂 RAG 架构}：最小闭环没跑通前，重排序、多查询、图检索只会把错误来源变得更难定位。
  \item \textbf{把所有文件一股脑入库}：低价值、历史版和重复文档会迅速拉高噪声，压垮最初的召回质量。
  \item \textbf{没有版本与权限元数据}：一旦答案引错来源或越权，团队无法知道问题是在检索器还是治理层。
  \item \textbf{\texttt{top-k} 固定过大}：材料一多，模型不一定更准，反而更容易忽略真正关键的条款。
  \item \textbf{只看最终答案，不看入模证据}：很多“模型答错”其实是在拼接前就已经丢了关键材料。
  \item \textbf{不设计拒答机制}：当系统必须看起来“总能回答”时，它几乎必然会在证据不足场景里编造。
\end{itemize}

再补四个在企业场景里尤其常见的取舍：
\begin{itemize}
  \item \textbf{先做覆盖面还是先做可靠性}：起步阶段通常先做可靠性。答对少数高价值问题，比模糊地答很多问题更有价值。
  \item \textbf{先做多数据源还是先做单一高价值源}：先把单一高价值源做稳，再扩源。数据源越多，治理成本越高。
  \item \textbf{先做更强模型还是先做更强日志}：对于最小闭环，日志通常比更强模型更重要，因为没有日志就不知道换模型是否真的解决了问题。
  \item \textbf{先提召回率还是先降噪声}：两者都重要，但企业条款问答更怕噪声过多导致答偏，因此通常要把“相关但无用”的片段治理好。
\end{itemize}

\section{本章练习}
\begin{exercise}[检索失败还是生成失败]
企业知识助手回答“VPN 权限申请只需要直属经理审批”，但日志显示召回结果里其实已经包含“安全负责人二次审批”的制度条款，且该条款也进入了最终提示。请判断这更像哪一层的问题，并说明第一步应该怎么修。
\end{exercise}

\begin{solution}
这更像{\heiti 生成层问题}。因为关键证据已经被召回，也进入了最终提示，说明检索层和拼接层至少没有先失手。第一步应该先收紧回答约束，例如要求模型必须显式列出审批链、必须引用对应条款，并在材料存在多条件时优先输出完整条件而不是只输出第一条结论。
\end{solution}

\begin{exercise}[为什么先做最小闭环]
为什么在企业知识助手里，不应该在最小 RAG 闭环还没跑稳时，就急着引入多查询、重排序、GraphRAG 或复杂 Agent 流程？
\end{exercise}

\begin{solution}
因为最小闭环没跑稳时，系统连“证据有没有被找到、有没有真正进入模型、模型是否忠于证据、知识是否还是最新版本”这几件基础事实都不能保证。此时叠加更多模块，只会增加延迟、成本和排障复杂度，却不一定解决真正的根因。先把最小闭环做成可检索、可引用、可拒答、可回放，后续优化才有明确抓手。
\end{solution}

\section{延伸阅读}
\begin{itemize}
  \item \textbf{Lewis et al., \emph{Retrieval-Augmented Generation for Knowledge-Intensive NLP Tasks}}：RAG 范式的经典起点，适合理解“检索先于生成”的基本思想。
  \item \textbf{LlamaIndex 或 LangChain 的 RAG 入门文档}：适合理解最小闭环在工程框架里通常如何分层实现。
  \item \textbf{Pinecone、Milvus、Weaviate 等检索系统官方文档}：适合理解索引、元数据过滤和召回接口在生产系统中的职责。
  \item \textbf{企业知识治理与文档生命周期资料}：适合理解 RAG 的上限不只由模型决定，也由知识源治理质量决定。
\end{itemize}

\section{小结}
RAG 基础章最重要的结论不是某个框架怎么用，而是企业知识助手必须先建立一条“可检索、可引用、可拒答、可回放”的最小闭环。这个闭环一旦建立，系统就第一次具备了从“聊天工具”迈向“知识系统”的条件；而这个闭环是否稳定，首先取决于文档准入、切分边界、元数据、日志和拒答规则这些看似朴素但最关键的基础工作。

\section{下一章过渡}
最小闭环建立后，新的瓶颈会很快出现：PDF 解析不稳、表格信息被切碎、用户问题过于口语化、召回结果“看着相关但并不好用”。下一章将集中处理这些文档理解与检索优化问题，回答企业知识助手应该按什么顺序把 RAG 做得更稳、更准。
